\section{The problem and an explicit
construction}\label{the-problem-and-an-explicit-construction}

A measurement of a vector \(x\in\mathbb C^8\) will mean a squared
modulus \(|q^*x|^2\), where \(q^*x=\sum_k\overline{q_k}x_k\). Every such
measurement is unchanged when \(x\) is multiplied by a complex number of
modulus one. The question is whether 26 choices of \(q\) can distinguish
every pair of vectors except for this unavoidable common phase. A family
with this property is called a phase-retrieving frame. We construct one
explicitly.

We identify a vector \(x\) with a polynomial \(f_x\) in a particular
eight-dimensional space. The measurement vectors will lie on a single
polynomially parametrized curve: their measurements are simply four
times the values \(|f_x(t)|^2\) at positive real parameters \(t\). The
construction has to achieve two different things. It must make 26
samples determine the entire polynomial \(|f_x(t)|^2\), and it must make
that polynomial determine \(f_x\) up to common phase.

Here is the argument in advance. The squared moduli belong to a
26-dimensional real polynomial space, because three coefficients always
vanish. Rolle's theorem then makes any 26 distinct positive samples
sufficient. To recover the polynomial itself from its squared modulus,
we align the first nonzero coefficients of two possible solutions and
compare their average \(v\) and half-difference \(h\). Equal squared
moduli say that \(v(t)\) and \(h(t)\) are perpendicular in the real
plane \(\mathbb C\) for every real \(t\). Their first two coefficients
leave only four cases. In three cases a positive sum of squares forces
the difference to vanish. The last case reduces to a real symmetric
two-by-two matrix whose diagonal entries have opposite signs; its
determinant can never be zero. The calculation below supplies the exact
entries of this matrix and covers all cases, including vanishing initial
coordinates.

For a polynomial \(f(t)=\sum c_k t^k\), write
\(f^\#(t)=\sum\overline{c_k}t^k\). Thus \(ff^\#\) is a real polynomial
and equals \(|f(t)|^2\) when \(t\) is real. We call it the
squared-modulus polynomial of \(f\).

The basis below is assembled from three short polynomials and a shifted
copy of them. Write \(a=(a_0,a_1,a_2)^T\) and set

\[
a=\begin{pmatrix}1+it+t^3+4it^5\\t^2+\tfrac i2t^3-t^5\\t^4+it^5\end{pmatrix},
\qquad
J=\begin{pmatrix}8&0&2\\0&1&0\\2&0&0\end{pmatrix}.
\]

The matrix \(J\) records a coefficient of the symmetric product:
\([t^5](aa^T)=iJ\), where \([t^5]\) denotes the coefficient of \(t^5\).
This identity explains the choice of \(J\); Section 2 shows how it
produces the required cancellation.

Take the three entries of \(a+t^8J^{-1}a^\#\), followed by
\(t^6-2it^7\), the three entries of \(t^8a\), and finally \(t^{14}\).
The conjugated correction is chosen to preserve the three missing
coefficients of the squared modulus. It also prevents the low-degree
part and its shifted copy from producing a simple ambiguity.

Indeed, without that correction the comparison space would contain both
\(a_0\) and \(t^8a_0\). The two nonproportional polynomials

\[
a_0(1+it^8),\qquad a_0(1-it^8)
\]

would then have identical squared moduli on the real line. The
correction breaks this particular construction of an ambiguity; the
proof below must still exclude all other pairs. The coefficients,
including the middle coefficient \(-2i\), are an explicit choice for
which the subsequent strict inequalities hold. Expanded, the eight basis
elements are

\[
\begin{aligned}
p_0(t)&=1+it+t^3+4it^5+\tfrac12t^{12}-\tfrac i2t^{13},\\
p_1(t)&=t^2+\tfrac i2t^3-t^5+t^{10}-\tfrac i2t^{11}-t^{13},\\
p_2(t)&=t^4+it^5+\tfrac12t^8-\tfrac i2t^9+\tfrac12t^{11}-2t^{12},\\
p_3(t)&=t^6-2it^7,\\
p_4(t)&=t^8+it^9+t^{11}+4it^{13},\\
p_5(t)&=t^{10}+\tfrac i2t^{11}-t^{13},\\
p_6(t)&=t^{12}+it^{13},\\
p_7(t)&=t^{14}.
\end{aligned}
\tag{1}
\]

Let \(t_1,\ldots,t_{26}\) be any distinct positive real numbers, and set

\[
q_j=2\bigl(\overline{p_0(t_j)},\ldots,\overline{p_7(t_j)}\bigr)^T
\in\mathbb C^8. \tag{2}
\]

\textbf{Theorem.} For all \(x,y\in\mathbb C^8\),

\[
|q_j^*x|^2=|q_j^*y|^2\quad(1\le j\le26)
\quad\Longrightarrow\quad y=e^{i\theta}x
\]

for some real \(\theta\). In particular, take \(t_j=j\); then all
entries of the 8-by-26 matrix \(Q=(q_1,\ldots,q_{26})\) are Gaussian
integers.

Every \(q_j\) is nonzero, since its last entry is \(2t_j^{14}\). Writing
\(m_{\mathbb C}(8)\) for the smallest possible number of such
measurements, the theorem gives \(m_{\mathbb C}(8)\le26\). It does not
assert that 25 measurements are impossible. Each measurement is given by
the single vector \(q_j\) in (2), or equivalently by the rank-one matrix
\(q_jq_j^*\). No more general quadratic measurements are used.

\section{Why 26 positive samples determine the squared
modulus}\label{why-26-positive-samples-determine-the-squared-modulus}

Put \(U=\operatorname{span}_{\mathbb C}(p_0,\ldots,p_7)\) and
\(f_x=\sum_{k=0}^7x_kp_k\). The orders of vanishing of these basis
polynomials at zero are \(0,2,\ldots,14\), and their first coefficients
are all 1. They are therefore linearly independent: in any nonzero
linear combination, the term of lowest order cannot cancel. In
particular \(x\mapsto f_x\) is an identification of \(\mathbb C^8\) with
\(U\).

A squared-modulus polynomial has degree at most 28 and would normally
have 29 real coefficients. The three linear restrictions below explain
the number 26 in this construction. Here \([t^k]\) means the coefficient
of \(t^k\).

For every \(u\in U\),

\[
[t^1](uu^\#)=[t^5](uu^\#)=[t^{13}](uu^\#)=0. \tag{3}
\]

These identities hold for all complex coefficients of \(u\). They can be
checked directly from (1), or by verifying that the three coefficients
vanish in all products \(p_jp_k^\#\). The exact checker does the latter,
checking real and imaginary parts separately. For example, writing
\(u=\sum z_jp_j\), its first six coefficients are

\[
z_0,\ iz_0,\ z_1,\ z_0+\tfrac i2z_1,\ z_2,
\ 4iz_0-z_1+iz_2.
\]

Consequently the coefficient of \(t^5\) in its squared-modulus
polynomial is

\[
2\operatorname{Re}\bigl(4i|z_0|^2-z_1\bar z_0+z_0\bar z_1
+\tfrac i2|z_1|^2\bigr)=0.
\]

To verify the degree-13 cancellation, return to the short vector \(a\)
and the real symmetric matrix \(J\) used to assemble (1). The corrected
low-degree entries are \(a+t^8J^{-1}a^\#\), the shifted entries are
\(t^8a\), and the remaining two are \(t^6-2it^7\) and \(t^{14}\). Direct
multiplication of the three short polynomials gives

\[
[t^5](a_i a_j^\#)=0,\qquad [t^5](a_i a_j)=iJ_{ij}.
\]

In products of two corrected low-degree entries, the coefficient of
degree 13 comes from the two cross terms. Their coefficient matrices are
\(iJJ^{-1}\) and \(-iJ^{-1}J\), whose sum is zero. A product of a
first-group entry and a conjugated tail entry has coefficient
\([t^5](a_i a_j^\#)=0\). Products involving the middle entry
\(t^6-2it^7\) have no degree-13 term: the low portions have degree at
most 12, the high portions start at degree 14, and the middle entry's
own norm is \(t^{12}+4t^{14}\). The remaining products start beyond
degree 13. This proves the three cancellations in (3) directly.

\textbf{Lemma 1 (positive sampling).} A nonzero real polynomial with at
most \(m\) nonzero monomials has at most \(m-1\) distinct positive
zeros.

\emph{Proof.} Induct on \(m\). Divide by its lowest power of \(t\),
which does not change positive zeros. The resulting polynomial has a
nonzero constant term and at most \(m-1\) other monomials. Its
derivative has at most \(m-1\) monomials and therefore at most \(m-2\)
distinct positive zeros by induction. Rolle's theorem proves the claim.
The case \(m=1\) is immediate. \(\square\)

The space of real polynomials of degree at most 28 with the three
coefficients in (3) equal to zero has dimension 26. The point of Lemma 1
is that evaluation at any 26 distinct positive numbers is injective on
this space; dimension counting alone would not establish that fact.

If the measurements of \(x,y\) agree, set \(f=\sum x_jp_j\) and
\(g=\sum y_jp_j\). By (2), \(q_j^*x=2f(t_j)\) and \(q_j^*y=2g(t_j)\), so
the measurements are four times the corresponding squared moduli. Their
difference of squared-modulus polynomials \(ff^\#-gg^\#\) has degree at
most 28 and misses the three exponents in (3). It therefore has at most
26 nonzero monomials. It vanishes at the 26 distinct positive sample
points, so Lemma 1 gives

\[
ff^\#=gg^\#. \tag{4}
\]

It remains to prove that (4) has only proportional solutions in \(U\).

\section{Comparing two polynomials with the same squared
modulus}\label{comparing-two-polynomials-with-the-same-squared-modulus}

We now prove that (4) determines a polynomial in \(U\) up to common
phase. The order of vanishing at zero organizes the comparison. The
nested spaces \(\operatorname{span}(p_i,\ldots,p_7)\) consist of the
elements of \(U\) vanishing to order at least \(2i\). Every nonzero
element belongs to exactly one of the differences between successive
spaces.

If either polynomial in (4) is zero, both are zero. Otherwise let \(i\)
be the first nonzero basis index of \(f\). Its squared-modulus
polynomial starts in degree \(4i\), with coefficient equal to the
squared modulus of that basis coefficient. The same must be true of
\(g\). Hence both have the same first index and leading coefficients of
equal modulus. Multiply one by a unit complex number to align these
coefficients, and divide both by the resulting common nonzero
coefficient. The squared-modulus identity is preserved, and the
coefficients at index \(i\) are now both 1.

Set

\[
v=\frac{f+g}{2},\qquad h=\frac{f-g}{2}.
\]

At every real \(t\), the exact identity
\(|v(t)+h(t)|^2-|v(t)-h(t)|^2=4\operatorname{Re}(h(t)\overline{v(t)})\)
shows what equal measurements require: the average and the difference
must be perpendicular in the real plane \(\mathbb C\) at every
parameter. This is an equality for the original pair of polynomials, not
a linear approximation to it. Since their first coefficients have been
aligned, \(h\) vanishes to higher order than \(v\). With all parameters
below real, we can write the following. Here \(x_j,y_j\) are new real
coordinates of the normalized average \(v\), not the coordinates of the
two original signals; \(r_j,a_j\) are the real and imaginary coordinates
of \(h\). These \(a_j\) are scalar coordinates, unrelated to the short
polynomials used to assemble the basis in Section 1.

\[
\begin{aligned}
v&=p_i+\sum_{j>i}(x_j+iy_j)p_j,\\
h&=\sum_{j>i}(r_j+ia_j)p_j,\\
F(t)&:=hv^\#+h^\#v=0.
\end{aligned}\tag{5}
\]

The index \(i\) ranges over all eight possibilities. In particular, we
have not assumed that the first coordinate of the original vector is
nonzero. Nor have we divided by a polynomial or by a measurement value.
Common factors, repeated zeros, and zero measurement values remain
within the argument.

The next coefficient distinguishes most of the remaining possibilities.
After its initial power of \(t\) is removed, each \(p_j\) starts at 1
and has a purely imaginary first derivative. Write this derivative as
\(is_j\). Thus

\[
p_j=t^{2j}(1+i s_jt)+O(t^{2j+2})
\]

where the real numbers \(s_j\) are

\[
(s_0,\ldots,s_7)=(1,\tfrac12,1,-2,1,\tfrac12,1,0). \tag{6}
\]

\textbf{Lemma 2 (comparison of the first two coefficients).} If
\(h\ne0\) in (5), and \(j\) is its first nonzero index, then \(r_j=0\)
and \(s_j=s_i\).

\emph{Proof.} The first two potentially nonzero coefficients of \(F\)
are \(2r_j\) and, after \(r_j=0\), \(2(s_i-s_j)a_j\). No larger basis
index can contribute to these two degrees. If the slopes differ, both
parts of the supposedly first nonzero coefficient vanish, a
contradiction. \(\square\)

The lemma says that the first possible nonzero difference has to have
the same first imaginary derivative as the leading basis element of the
average. From (6) the complete list is:

{\def\LTcaptype{none} % do not increment counter
\begin{longtable}[]{@{}ll@{}}
\toprule\noalign{}
first signal index \(i\) & possible first difference indices \\
\midrule\noalign{}
\endhead
\bottomrule\noalign{}
\endlastfoot
0 & 2, 4, 6 \\
1 & 5 \\
2 & 4, 6 \\
3 & none \\
4 & 6 \\
5, 6, 7 & none \\
\end{longtable}
}

For \(i=3,5,6,7\), no nonzero \(h\) is possible. It remains to consider
\(i=0,1,2,4\). The next section gives every equation needed for these
four cases.

\section{The coefficient equations force the difference to
vanish}\label{the-coefficient-equations-force-the-difference-to-vanish}

For each first index \(i\), put

\[
W_i=\operatorname{span}_{\mathbb C}(p_{i+1},\ldots,p_7),\qquad W_7=\{0\}.
\]

After normalization we have \(v\in p_i+W_i\) and \(h\in W_i\). For each
fixed \(v\), consider the real-linear map

\[
L_v:W_i\longrightarrow\mathbb R[t]_{\le28},\qquad
L_v(h)=hv^\#+h^\#v,
\]

where the complex vector space \(W_i\) is regarded as a real vector
space. Our remaining task is to prove \(\ker L_v=\{0\}\) for every such
\(v\). Pointwise perpendicularity alone would not imply this: for
example, \(h=icv\) with \(c\in\mathbb R\) satisfies \(hv^\#+h^\#v=0\).
The normalization excludes this nonzero constant-phase direction because
\(h\) has no \(p_i\) coefficient whereas \(v\) has coefficient 1. The
rest of the argument uses the requirement that \(h\) belongs to the
fixed polynomial space \(W_i\).

The low-order coefficients of \(L_v(h)\) determine most coordinates of
\(h\) successively; a few higher-order coefficients rule out the
remaining coordinates. Thus all calculations below concern the kernel of
one linear map at a time, for an arbitrary normalized average \(v\).

Write \(E_k=[t^k]F\). To extract any of these equations directly, write
\(v(t)=\sum_\ell v_\ell t^\ell\) and \(h(t)=\sum_\ell h_\ell t^\ell\);
these are ordinary monomial coefficients, with coefficients outside the
degree range understood to be zero. Then

\[
E_k=2\operatorname{Re}\sum_{\ell=0}^k
h_\ell\overline{v_{k-\ell}}.
\]

After the preceding variables have been determined, the equation
\(E_{2i+2j}=0\) has coefficient 2 on \(r_j\) and determines it uniquely.
The next equation has coefficient \(2(s_i-s_j)\) on \(a_j\). If this
number is nonzero, it determines \(a_j\) as well. If it is zero, we
retain \(a_j\) and keep any resulting equation as a necessary condition.
These operations apply to every actual solution of \(F=0\).

Every division is by a fixed nonzero rational number, never by a
variable expression that might vanish for a special signal. Thus no
exceptional cases are lost. All identities below follow by polynomial
multiplication and these substitutions. The exact checker reconstructs
them from (1), but the required equations are displayed here.

The order of the four cases puts the shortest calculations first. Most
end with a nonzero difference coefficient times a strictly positive sum
of squares. In the last case two coefficients remain together; their two
equations form a symmetric matrix with a determinant of fixed sign.

\subsection{\texorpdfstring{The average starts with
\(p_4\)}{The average starts with p\_4}}\label{the-average-starts-with-p_4}

By Lemma 2, a nonzero difference can first occur only at index 6. In
this case write

\[
\begin{aligned}
v&=p_4+(x_5+iy_5)p_5+(x_6+iy_6)p_6+(x_7+iy_7)p_7,\\
h&=(r_6+ia_6)p_6+(r_7+ia_7)p_7.
\end{aligned}
\]

Here is the short elimination in full. The first equation is
\(E_{20}=2r_6=0\). With \(r_6=0\), the coefficient \(E_{21}\) vanishes
identically because the relevant slopes agree, and the next two are

\[
E_{22}=2r_7+2a_6y_5,\qquad E_{23}=2a_7-a_6x_5.
\]

Setting these equal to zero gives

\[
r_6=0,\qquad r_7=-a_6y_5,\qquad a_7=\tfrac12a_6x_5.
\]

After these substitutions,

\[
E_{25}=\frac{a_6}{2}\left(x_5^2+2(y_5-2)^2+8\right). \tag{7}
\]

The factor in parentheses is strictly positive. Thus \(a_6=0\), and all
coefficients of \(h\) vanish.

\subsection{\texorpdfstring{The average starts with
\(p_1\)}{The average starts with p\_1}}\label{the-average-starts-with-p_1}

Only differences starting at 5 need consideration. Eliminating the
coefficients of degrees 12 through 17 gives

\[
r_5=0,\quad r_6=-a_5y_2,\quad a_6=a_5x_2,
\quad r_7=-a_5y_3,\quad a_7=5a_5x_3.
\]

The following coefficient is then a particularly useful necessary
constraint:

\[
E_{21}=-2a_5\left(1+x_2^2+y_2^2+10x_3^2+2y_3^2\right). \tag{8}
\]

It forces \(a_5=0\), hence \(h=0\). Unused coefficients are not presumed
to vanish automatically: it is sufficient that every actual solution
must satisfy (8).

\subsection{\texorpdfstring{The average starts with
\(p_2\)}{The average starts with p\_2}}\label{the-average-starts-with-p_2}

Start with differences at 4 or later. Processing degrees 12 through 16
gives

\[
\begin{aligned}
r_4&=0,\\
r_5&=-a_4y_3,\\
a_5&=6a_4x_3,\\
r_6&=a_4\left(-1-y_4+\tfrac52y_3-5x_3y_3\right).
\end{aligned}
\]

The degree-17 equation is

\[
E_{17}=-a_4\left(30x_3^2+5(y_3-\tfrac25)^2+\tfrac{46}{5}\right). \tag{9}
\]

Thus \(a_4=0\), and the coefficients of \(h\) at 4 and 5 are zero.
Restarting with a possible difference at 6, degrees 16 through 19 give

\[
r_6=0,\qquad r_7=-a_6y_3,\qquad a_7=3a_6x_3.
\]

Now

\[
E_{21}=-2a_6(1+6x_3^2+2y_3^2). \tag{10}
\]

Hence \(a_6=0\), and again \(h=0\).

\subsection{\texorpdfstring{The average starts with \(p_0\): the first
possible
difference}{The average starts with p\_0: the first possible difference}}\label{the-average-starts-with-p_0-the-first-possible-difference}

Process degrees 2 through 8. The solved variables are

\[
\begin{aligned}
r_1=a_1=r_2&=0,\\
r_3&=-a_2y_1,\qquad a_3=\tfrac16a_2x_1,\\
r_4&=a_2\left(1-y_2-\tfrac52y_1+\tfrac56x_1y_1\right).
\end{aligned}
\]

The degree-9 coefficient becomes

\[
E_9=a_2\left(\tfrac56x_1^2+5(y_1-\tfrac25)^2+\tfrac{46}{5}\right). \tag{11}
\]

It follows that \(a_2=0\). Consequently every coefficient of \(h\)
before index 4 is zero.

\subsection{\texorpdfstring{The average starts with \(p_0\): the last
two
coefficients}{The average starts with p\_0: the last two coefficients}}\label{the-average-starts-with-p_0-the-last-two-coefficients}

Now start \(h\) at index 4. The coefficients of degrees 8 through 15
give

\[
\begin{aligned}
r_4&=0,& r_5&=-a_4y_1,&a_5&=a_4x_1,\\
r_6&=-a_4y_2,\\
r_7&=-a_4y_3-a_6y_1+a_4x_2y_1,\\
a_7&=3a_4x_3+\tfrac12a_6x_1-\tfrac12a_4x_1x_2.
\end{aligned}\tag{12}
\]

There are two remaining independent imaginary difference variables,
\(a_4,a_6\). Introduce the real polynomials

\[
\begin{aligned}
R&=x_1^2+2(y_1-2)^2+8,\\
T&=-2x_1x_3-4y_1y_3-2x_2,\\
S&=-1-x_1^2-y_1^2-6x_3^2-2y_3^2-x_2^2-y_2^2,\\
\rho&=a_6-a_4x_2.
\end{aligned}\tag{13}
\]

After (12), direct coefficient extraction gives

\[
E_{17}=\frac{R\rho+2a_4T}{2},\qquad
E_{21}=2a_4S+T\rho. \tag{14}
\]

The last two equations have a simple two-dimensional meaning. In the
coordinates \((2a_4,\rho)\) for the remaining difference, (14) reads

\[
\begin{pmatrix} S&T\\ T&R\end{pmatrix}
\begin{pmatrix}2a_4\\ \rho\end{pmatrix}
=\begin{pmatrix}E_{21}\\2E_{17}\end{pmatrix}.
\]

The two diagonal entries have fixed opposite signs: \(S\le-1\) and
\(R\ge8\). The associated quadratic form is therefore negative on one
coordinate axis and positive on the other. In this two-dimensional
system the off-diagonal entry \(T\), whatever its sign, contributes
\(-T^2\) to the determinant. Hence \(RS-T^2\le-8\), and the matrix is
invertible for every average \(v\).

The same argument can be written as a scalar elimination identity, with
no division by the determinant:

\[
\boxed{R E_{21}-2T E_{17}=2a_4(RS-T^2).}\tag{15}
\]

For a pair with equal squared moduli, \(E_{17}=E_{21}=0\), so (15)
forces \(a_4=0\). Multiplying the same identity by \(-a_4\) gives its
quantitative form:

\[
-a_4(RE_{21}-2TE_{17})=2a_4^2(R(-S)+T^2)\ge16a_4^2.
\]

This inequality applies to either sign of \(a_4\). With \(a_4=0\), the
first equation in (14) gives \(a_6=0\), and (12) makes every remaining
coefficient of \(h\) zero.

Both rows of the two-by-two system are coefficients of the same
polynomial \(F\), for the same average \(v\) and difference \(h\). Thus
their joint use applies to every candidate pair in (4). The remaining
coefficients of \(F\) need not be solved: every actual solution must
satisfy the equations already used, and those equations force \(h=0\).

\section{Conclusion and the equivalent matrix
statement}\label{conclusion-and-the-equivalent-matrix-statement}

All eight possible first signal indices have been covered. In every case
\(h=0\), so after normalization \(f=g\). Undoing that normalization
gives precisely a common unit phase. Section 2 then proves the
finite-measurement theorem.

For completeness, the conclusion can also be expressed entirely in terms
of Hermitian matrices. Define the real-linear map
\(\mathcal A_Q(H)=(q_j^*Hq_j)_{j=1}^{26}\) on the Hermitian
\(8\)-by-\(8\) matrices. Equality of the measurements of \(x\) and \(y\)
is exactly the equation \(\mathcal A_Q(xx^*-yy^*)=0\).

The columns of \(Q\) span \(\mathbb C^8\): a vector perpendicular to
them would give a polynomial in \(U\) vanishing at 26 distinct points,
hence the zero polynomial since its degree is at most 14. If a positive
semidefinite Hermitian matrix \(H\) belongs to the kernel of
\(\mathcal A_Q\), each \(q_j^*Hq_j=0\) implies \(Hq_j=0\). Spanning then
gives \(H=0\). The same applies after changing the sign of a negative
semidefinite matrix. Finally, every indefinite Hermitian matrix of rank
two can be written as \(xx^*-yy^*\) with \(x,y\) not proportional, by
its spectral decomposition. The theorem has excluded precisely these
remaining possibilities. Therefore

\[
\ker\mathcal A_Q\cap\{H=H^*:0<\operatorname{rank}H\le2\}=\varnothing.
\]

\begingroup\small

\section{Scope and verification}\label{scope-and-verification}

The proof uses only the displayed polynomials, coefficient conjugation,
elementary real inequalities, the normalization by the first nonzero
coefficient, and Rolle's theorem. It does not depend on any prior
unpublished classification or obstruction claim. The construction was
motivated by choosing \(t^{14}\), a polynomial with a repeated root, as
an element shared by \(U\) and \(U^\#\), where \(U^\#=\{f^\#:f\in U\}\).
No theory of such spaces is needed here.

The result concerns the specified frames and gives
\(m_{\mathbb C}(8)\le26\). It asserts neither optimality nor success for
every or generic 26-vector frame. Numerical conditioning and a practical
reconstruction algorithm with noise guarantees have not been
established.

The supplementary archive provides the exact Gaussian-integer matrix
\texttt{PR26\_FRAME\_QI.json}, the paper's sources, both verification
implementations, and their logs. The original standard-library
implementation uses exact rational arithmetic. An isolated rerun passed
all 336 polynomial checks, seven matrix checks, and two negative
controls, with mathematical results matching the supplied logs. These
checks include the three missing coefficients, all substitutions and
case identities, the actual frame entries, and rank calculations modulo
a prime.

A separate SymPy implementation transcribes (1) without importing code
from the original verifier. Its 633 finite checks passed: all 64 complex
products, all 28 possible first-index pairs, the substitutions by fixed
nonzero coefficients, the strict-sign identities, and all 208 matrix
entries. The newly displayed symmetric two-by-two formulation is just a
rearrangement of (14); its determinant agrees with (15). A separately
supplied Wolfram log reports agreement for the three missing
coefficients and eight principal coefficient identities and their
combinations; that engine was not rerun during the local review.

The sampling argument, normalization, and case coverage have received a
local logical review. Computational check counts do not count
independent mathematical proofs. This is a proof candidate for critical
examination: external independent expert review and proof-assistant
verification have not been completed.

The mathematics, exploration, and verification were developed with
substantial AI assistance. Detailed research records are retained
separately. The complete proposed proof is given here, and the exact
input and reproducibility scripts accompany it.

\endgroup
